%% kapa: exam=1 round=37 year=2026 period=2 form=A subject=appraisal_law subject_ko=감정평가관계법규 seq=10 number=10
%% source: Q-Net 감정평가사 기출문제 (cst00309) artlSeq=5258924 file=raw/1008/5258924/2026년 제37회 감정평가사 1차 2교시.pdf
%% conversion: kordoc 4.24.51 -> tools/gen1.py (auto); review=proofread-sonnet; flags=none
\begin{question}{10}
국토의 계획 및 이용에 관한 법령상 도시ㆍ군계획시설사업의 시행에 관한 설명으로 옳지 않은 것은?
\begin{choices}
  \item 도시ㆍ군계획시설결정의 고시일부터 20년이 되기 전에 실시계획이 효력을 잃고 다른 도시ㆍ군계획시설사업이 시행되지 아니하는 경우, 그 결정은 도시ㆍ군계획시설결정의 고시일부터 20년이 되는 날의 다음 날에 효력을 잃는다.
  \item 도시ㆍ군계획시설사업을 분할시행하는 때에는 분할된 지역별로 실시계획을 작성할 수 있다.
  \item 행정청이 아닌 도시ㆍ군계획시설사업의 시행자는 이해관계인에게 공시송달을 하려는 경우에는 국토교통부장관, 관할 시ㆍ도지사 또는 대도시 시장의 승인을 받아야 한다.
  \item ｢빈집 및 소규모주택 정비에 관한 특례법｣에 따라 도시ㆍ군관리계획의 결정이 의제되는 경우에는 해당 도시ㆍ군계획시설결정의 고시일부터 1년 이내에 단계별 집행계획을 수립하여야 한다.
  \item 시ㆍ도지사나 대도시 시장은 공사완료보고서를 받으면 지체 없이 준공검사를 하여야 한다.
\end{choices}
\end{question}
\answer{4}
